\section{Complex Engineering Problem}

This project is a complex engineering problem. It is not just writing code. We need to understand math, build numerical solvers, design neural networks, and compare results carefully. The Lorenz-1960 system is nonlinear and coupled, so solving it with an ANN requires real engineering decisions, not just running a script.

\subsection{Complex Problem Solving}

Table~\ref{tab:p_solve} shows how this project maps to the complex problem solving categories.

\begin{center}
    \begin{table}[ht]
        \centering
        \caption{Mapping with complex problem solving.}
        \begin{tabular}{|p{0.12\textwidth}|p{0.12\textwidth}|p{0.12\textwidth}|p{0.12\textwidth}|p{0.12\textwidth}|p{0.12\textwidth}|p{0.12\textwidth}|}
        \hline
        P1& P2& P3& P4& P5& P6& P7\\
        Dept of Knowledge & Range of Conflicting Requirements & Depth of Analysis & Familiarity of Issues & Extent of Applicable Codes & Extent of Stakeholder Involvement & Inter-dependence\\
        \hline 
        $\checkmark$ & $\checkmark$ & $\checkmark$ & $\times$ & $\times$ & $\times$ & $\checkmark$\\
        &&&&&&\\
        \hline 
        \end{tabular}
        \label{tab:p_solve}
    \end{table}
\end{center}

\subsubsection{P1: Depth of Knowledge --- Applicable ($\checkmark$)}

We need to know many things at once: how ODEs work, how numerical solvers like Runge-Kutta produce reference solutions, how neural networks learn, how optimizers like Adam update weights, and how to measure errors. Knowing just one of these is not enough. They all connect together.

\subsubsection{P2: Range of Conflicting Requirements --- Applicable ($\checkmark$)}

We face trade-offs. A bigger network may give better accuracy but takes longer to train and may overfit. A faster training setup may be unstable. A model that works great once may not work again with a different random seed. We cannot get the best of everything at the same time, so we have to find a good balance.

\subsubsection{P3: Depth of Analysis --- Applicable ($\checkmark$)}

We do not just run the code and check if it works. We compare many ANN setups, look at error numbers, study how training behaves, check if results are stable, and explain why one setup works better than another. That takes real analysis.

\subsubsection{P4: Familiarity of Issues --- Not Applicable ($\times$)}

The Lorenz-1960 system is a well-known benchmark, and the techniques we use---Runge-Kutta solvers, feedforward ANNs, Adam optimizer---are all established methods. Although our specific combination is novel, the underlying issues are not unfamiliar to practitioners. So this criterion does not apply.

\subsubsection{P5: Extent of Applicable Codes --- Not Applicable ($\times$)}

Our project is a computational research study, not a commercial or safety-critical system. There are no industry standards, regulations, or professional codes that constrain how we design the ANN or evaluate its accuracy. So this criterion does not apply.

\subsubsection{P6: Extent of Stakeholder Involvement --- Not Applicable ($\times$)}

The only stakeholders are the student researchers and the academic supervisors. There are no external clients, end-users, or regulatory bodies whose conflicting needs must be balanced. All decisions are based on scientific reasoning, not stakeholder demands. So this criterion does not apply.

\subsubsection{P7: Interdependence --- Applicable ($\checkmark$)}

Every part of this project depends on the other parts. The solver makes the training data. The data affects how the ANN learns. The ANN design affects the results. The optimizer affects training. If any one part goes wrong, the whole result is affected.


\subsection{Engineering Activities}

Table~\ref{tab:e_act} shows how this project maps to the complex engineering activity categories.

\begin{center}
    \begin{table}[ht]
        \centering
        \caption{Mapping with complex engineering activities.}
        \begin{tabular}{|p{0.18\textwidth}|p{0.18\textwidth}|p{0.18\textwidth}|p{0.18\textwidth}|p{0.18\textwidth}|}
        \hline
        A1& A2& A3& A4& A5\\
        Range of resources & Level of Interaction & Innovation & Consequences for society and environment & Familiarity\\
        \hline 
        $\checkmark$ & $\checkmark$ & $\checkmark$ & $\times$ & $\checkmark$\\
        &&&&\\
        \hline 
        \end{tabular}
        \label{tab:e_act}
    \end{table}
\end{center}

\subsubsection{A1: Range of Resources --- Applicable ($\checkmark$)}

We use many different tools and knowledge sources: ODE math, Runge-Kutta solvers, SciPy, PyTorch, NumPy, Matplotlib, and 13 research papers. No single tool or source is enough to complete this project.

\subsubsection{A2: Level of Interaction --- Applicable ($\checkmark$)}

The parts of this project talk to each other. The solver gives data to the ANN. The ANN design affects training. The training setup affects results. The evaluation depends on everything before it. You cannot work on one part without thinking about the others.

\subsubsection{A3: Innovation --- Applicable ($\checkmark$)}

No one has done a systematic comparison of plain ANN architectures on the Lorenz-1960 system before. Most existing work uses PINNs or hybrid methods \cite{matthews2024pinnde}. We test a different idea: can a plain ANN, without physics in the loss function, still solve this problem well? That is the new part.

\subsubsection{A4: Consequences for Society and Environment --- Not Applicable ($\times$)}

Our project is a purely academic study that trains small neural networks on a personal computer. It does not produce any physical product, does not affect public health or safety, and uses minimal computational resources. The Lorenz-1960 system is a theoretical benchmark, so our results do not drive any real-world decisions. This criterion does not apply.

\subsubsection{A5: Familiarity --- Applicable ($\checkmark$)}

This is not a standard textbook problem. The Lorenz-1960 system is nonlinear and chaotic. There is no ready-made ANN setup that we can just use. We have to figure out the best architecture through experiments, and the results need careful interpretation because low training loss does not always mean the solution is correct everywhere.
